\documentclass[11pt]{article}
\usepackage[margin=1in]{geometry}
\usepackage{amsmath,amssymb,amsthm,booktabs,hyperref}
\newtheorem{theorem}{Theorem}
\newtheorem{lemma}{Lemma}
\title{Extremal rows in Castronuovo's matrix construction}
\author{Principia Math In Progress}
\date{28 September 2026}
\begin{document}
\maketitle
\begin{abstract}
We prove the extremal-row conjecture, Conjecture 6.5 of Castronuovo
\cite{C}, by determining the entire coordinatewise maximal frontier
of each positive generation. This gives the two claimed integer
sequences and the exact multiplicities in their extremal rows.
The initial seed and the offset of the odd sequence are stated explicitly.
\end{abstract}

\section*{Task 1: the original problem and its precise indexing}
In row--column coordinates the two construction maps are
\[
 G(n,k)=(n+k,k),\qquad S(n,k)=(3n+1-2k,2n+1-k).
\]
The positive tree begins at $(1,1)$ in Step 1. Let $C_i$ be its
generation at Step $i$, so $C_1=\{(1,1)\}$ and
$C_{i+1}=G(C_i)\cup S(C_i)$. Put
$r_i=\max\{n:(n,k)\in C_i\}$ and set $r_0=0$ for the initial
seed $(0,0)$. This is the positive-orbit interpretation used in
Sections 2 and 3 of \cite{C}.

There are two harmless indexing conventions to make explicit.
The current OEIS sequence A061278 begins $0,1,5,20,76,\ldots$;
the odd rows begin $1,5,20,76,\ldots$. Also, the multiplicity-two
assertion concerns positive even steps: Step 0 consists of one seed.
These conventions do not supply a counterexample to the intended
question. We prove the positive-step assertions in full.

\begin{theorem}\label{thm:main}
For $j\geq0$,
\[
 r_{2j}=\mathrm{A001571}(j),\qquad
 r_{2j+1}=\mathrm{A061278}(j+1),
\]
with the OEIS sequences indexed from zero. Equivalently,
\[
 r_0=0,\quad r_1=1,\quad r_2=2,\quad r_3=5,\qquad
 r_{i+4}=4r_{i+2}-r_i+1\quad(i\geq0).
\]
For every $i\geq1$, exactly two distinct elements of $C_i$ lie in
row $r_i$ when $i$ is even, and exactly one when $i$ is odd.
\end{theorem}

\section*{A two-point frontier}
Change coordinates by
\[
 a=2k,\qquad b=2n+1-2k,\qquad n=\frac{a+b-1}{2}.
\]
The seed is $(a,b)=(2,1)$, and the maps become
\[
 G(a,b)=(a,a+b),\qquad S(a,b)=(a+2b,b).
\]
Both coordinates stay positive. We say $x$ strictly dominates $y$
when both coordinates of $x$ are at least those of $y$ and at least
one inequality is strict. Both shears preserve strict domination:
they have nonnegative entries and are injective. A strictly
dominated point has strictly smaller coordinate sum.

The two image cones are disjoint: a $G$ image has $b>a$, whereas an
$S$ image has $a>2b$. Moreover, the seed lies in neither cone.
Reading the final letter and applying its inverse therefore shows
that two different words, including words of different lengths,
cannot have the same image of the seed. Consequently the tree
indeed enumerates distinct entries, with $|C_i|=2^{i-1}$.

\begin{lemma}\label{lem:frontier}
At Step $2j$, $j\geq1$, the coordinatewise maximal points are precisely
\[
 A=(u,v),\qquad B=(2v-u,2u-v),\qquad 0<u<v<2u,
\]
where
\[
 \binom{u}{v}=
 \begin{pmatrix}1&2\\1&3\end{pmatrix}^{j-1}\binom{2}{3}.
\]
Every other point is strictly dominated by one of these two points.
At Step $2j+1$ the maximal points are precisely
\[
 D=(2v-u,u+v),\qquad E=(u+2v,v),
\]
and all other points are strictly dominated by one of them.
\end{lemma}

\begin{proof}
Step 2 consists of $(2,3)$ and $(4,1)$, proving the initial case.
Suppose the asserted even frontier is known. Its four children are
\[
 GA=(u,u+v),\quad GB=(2v-u,u+v),\quad
 SA=(u+2v,v),\quad SB=(3u,2u-v).
\]
The point $GB$ strictly dominates $GA$, since their first-coordinate
difference is $2(v-u)>0$. The point $SA$ strictly dominates $SB$,
since both coordinate differences are $2(v-u)>0$.
Thus the surviving points are $D=GB$ and $E=SA$.
They are incomparable: $E$ has the larger first coordinate and
$D$ the larger second coordinate. Children of previously dominated
points remain strictly dominated, so no other maximal points occur.

The four children of $D,E$ at the next even step are
\[
 GD=(2v-u,3v),\quad SD=(u+4v,u+v),\quad
 GE=(u+2v,u+3v),\quad SE=(u+4v,v).
\]
Here $GE$ strictly dominates $GD$, with differences $(2u,u)$;
$SD$ strictly dominates $SE$, with differences $(0,u)$.
The remaining pair $GE,SD$ is incomparable. Put
$u'=u+2v$ and $v'=u+3v$. Then
\[
 GE=(u',v'),\qquad SD=(2v'-u',2u'-v'),
\]
and $0<u'<v'<2u'$ because $v'-u'=v>0$ and
$2u'-v'=u+v>0$. Strict domination again discards all other
children. This proves the induction and the matrix formula.
\end{proof}

\section*{Rows, multiplicities, and the recurrence}
At even Step $2j$ the two frontier sums are both $u+v$.
Every other point has a smaller sum. Hence
\[
 r_{2j}=\frac{u+v-1}{2},\qquad
 \text{maximizing columns }\frac u2,\ \frac{2v-u}{2}.
\]
At the following odd step the two frontier sums are $3v$ and
$u+3v$. Since $u>0$, only $E$ maximizes the sum, and
\[
 r_{2j+1}=\frac{u+3v-1}{2},\qquad
 \text{maximizing column }\frac{u+2v}{2}.
\]
The two even points are distinct since $v>u$; all these column
indices are integral since $u$ is even and $v$ odd, as follows
inductively from the update. Step 1 has a single point, as required.

The update matrix $P$ satisfies $P^2-4P+I=0$. Therefore every
linear form in its iterates satisfies $z_{j+2}=4z_{j+1}-z_j$.
Subtracting 1 and dividing by 2 in the displayed row formulas gives
the asserted inhomogeneous recurrence for each parity, initially
for $i\geq2$. Directly, $r_4=9=4r_2-r_0+1$ and
$r_5=20=4r_3-r_1+1$, which supplies the two missing cases.
The OEIS entries \cite{even,odd} give this same recurrence with
initial values $(0,2)$ and $(0,1)$ respectively. Uniqueness of
recurrence solutions proves the sequence identifications in
Theorem~\ref{thm:main}.

For example, at Step 6, $(u,v)=(30,41)$, so the maximal row is 35
and its two new entries have columns 15 and 26. This recovers
Example 6.6 in \cite{C}.

\paragraph{The origin self-loop.}
If the introductory rule in \cite{C} is applied literally to the
origin at every step, then $G(0,0)=(0,0)$ repeatedly propagates old
positions. The set of positions changed at Step $i$ is then
$\{(0,0)\}\cup C_1\cup\cdots\cup C_i$.
This does not change the extremal-row answer: on the positive tree
$1\leq k\leq n$, and both maps strictly increase the row
($n+k>n$ and $3n+1-2k>n$). Thus $r_{i+1}>r_i$ and none of
the older generations reaches the new maximum. The same maxima
and multiplicities hold under that reading as well.

\section*{Proposed follow-up Tasks 2--5}
These are further questions, not claims established here.
\begin{enumerate}
\setcounter{enumi}{1}
\item For an arbitrary positive integer seed $(a_0,b_0)$, classify
the coordinatewise maximal frontiers under $G(a,b)=(a,a+b)$ and
$S(a,b)=(a+2b,b)$, and determine the maximal coordinate sum at
each word length.
\item Replace $S$ by $S_d(a,b)=(a+db,b)$ for an integer $d\geq1$.
Starting from $(2,1)$, determine for which $d$ the maximal sums
obey a linear recurrence and determine their multiplicities.
\item For the original construction, determine the second-largest
distinct row in every positive generation and the number of entries
of that generation in that row, including all small-step exceptions.
\item Choose the letters $G,S$ independently with probability $1/2$.
For the resulting row $N_i$ at Step $i$, prove or disprove a central
limit theorem for $\log N_i$ and, if it holds, characterize its drift
and asymptotic variance.
\end{enumerate}

\paragraph{Verification and status.}
The accompanying exact program enumerates the original row--column
tree through Step 18 and checks every frontier, extremal row, column,
and multiplicity against the formulas. These finite checks supplement
the proof. Current-source and bounded public-claim checks are recorded
separately. Review is by the authoring agent, not independent peer review.

\begin{thebibliography}{3}
\bibitem{C} N. Castronuovo, \emph{On the number of fixed points of the
map $\gamma$}, arXiv:2102.02739v4, 25 October 2023,
Conjecture 6.5 and Example 6.6, pp. 13--14.
\url{https://arxiv.org/abs/2102.02739}.
\bibitem{even} OEIS Foundation, \emph{A001571}, accessed 28 September 2026.
\url{https://oeis.org/A001571}.
\bibitem{odd} OEIS Foundation, \emph{A061278}, accessed 28 September 2026.
\url{https://oeis.org/A061278}.
\end{thebibliography}
\end{document}
